\documentclass[12pt]{report}

\usepackage{graphicx}
\usepackage{amsmath}
\usepackage{booktabs}
\usepackage{geometry}
\usepackage{xcolor}
\usepackage{pagecolor}
\geometry{margin=1in}

\begin{document}

% ================= TITLE PAGE =================
\begin{titlepage}
\pagecolor{white}
\color{black}

\centering

\vspace*{1cm}

{\Huge \textbf{Real-Time Road Mapping Parallel Image Processing Pipeline}}\\[1.5cm]

{\large A Project Component for}\\[0.2cm]
{\large \textbf{Parallel and Distributing Computing (UCS645)}}\\[1.5cm]

{\large By}\\[0.3cm]
{\large \textbf{SAMIA SHARIF}}\\
{\large 102303840}\\[1.5cm]

{\large Under the guidance of}\\[0.3cm]
{\large Dr. Saif Nalband}\\
{\small (Assistant Professor, DCSE)}\\[1.5cm]

\includegraphics[width=0.4\textwidth]{tietlogo.jpg}\\[1.5cm]

{\small DEPARTMENT OF COMPUTER SCIENCE AND ENGINEERING}\\
{\small THAPAR INSTITUTE OF ENGINEERING AND TECHNOLOGY}\\
{\small (DEEMED TO BE UNIVERSITY)}\\
{\small PATIALA - 147004}\\[0.5cm]

{\small APRIL-MAY, 2026}

\end{titlepage}

\nopagecolor
\newpage

% ================= CONTENT =================
\tableofcontents
\newpage

% ---------------- INTRODUCTION ----------------
\chapter{Introduction}

\section{Background}
Modern autonomous and intelligent systems require processing of large volumes of data in real-time. Traditional sequential systems are inefficient due to high latency and inability to utilize computational resources effectively.

\section{Introduction to Problem Statement}
This project explores the application of pipeline parallelism in real-time image processing systems using Python-based multithreading. Additionally, OpenMP is utilized in a separate C++ module to demonstrate shared-memory parallelism and data-level parallel computation. The combination of these approaches provides a comprehensive understanding of different parallel computing paradigms.

% ---------------- LITERATURE ----------------
\chapter{Literature Review}

Parallel computing improves efficiency through pipeline parallelism and asynchronous processing.
\renewcommand{\arraystretch}{1.3}
\begin{table}[h]
\centering
\begin{tabular}{|l|l|l|l|l|}
\hline
OpenMP & Shared-memory parallelism & Parallel for, sections, reduction & Efficient CPU parallelism & Limited to C/C++ \\ \hline
Reference & Concept & Technique & Significance & Limitations \\ \hline
YOLO & Detection & CNN & Fast detection & No system design \\ \hline
OpenCV & Image Processing & Edge Detection & Feature extraction & No parallelism \\ \hline
\end{tabular}
\caption{Summary of Literature Review}
\end{table}

% ---------------- RESEARCH GAP ----------------
\chapter{Research Gaps}
\begin{itemize}
\item Lack of pipeline-based systems
\item Focus only on algorithms
\end{itemize}

% ---------------- PROBLEM ----------------
\chapter{Problem Formulation}

\[
T_{seq} = T_1 + T_2 + T_3 + T_4
\]

\[
T_{parallel} \approx \max(T_1, T_2, T_3, T_4)
\]

% ---------------- OBJECTIVES ----------------
\chapter{Objectives}
\begin{itemize}
\item Design a multi-stage processing pipeline
\item Implement parallel execution using Python threads
\item Use queues for inter-stage communication
\item Implement OpenMP-based parallel computation in C++
\item Compare sequential, pipeline, and OpenMP performance
\item Demonstrate real-time system behavior
\end{itemize}

% ---------------- METHODOLOGY ----------------
\chapter{Methodology}
\section{OpenMP-Based Parallel Implementation}

To complement the pipeline-based parallel system, an OpenMP-based implementation is developed in C++ to demonstrate shared-memory parallelism.

OpenMP is a widely used API for parallel programming in C/C++ that allows efficient utilization of multi-core processors.

\subsection*{Implementation Details}
\begin{itemize}
\item Parallel Sections: Used to simulate pipeline stages
\item Parallel For Loops: Used for data parallel operations
\item Reduction: Used for aggregation operations
\item Dynamic Scheduling: Improves load balancing
\end{itemize}

\subsection*{Equation}
\[
T_{openmp} \approx \frac{T_{seq}}{N}
\]
\section*{Pipeline}
Input $\rightarrow$ Preprocess $\rightarrow$ Detect $\rightarrow$ Lane $\rightarrow$ Output

\section*{Algorithm}
\begin{verbatim}
1. Read image
2. Preprocess
3. Detect objects
4. Detect lanes
5. Display output
\end{verbatim}

\section{Equation}

\[
Speedup = \frac{T_{seq}}{T_{parallel}}
\]

\[
Efficiency = \frac{Speedup}{N}
\]

% ---------------- DATASET ----------------
\chapter{Datasets}

BDD100K dataset is used for testing.
\renewcommand{\arraystretch}{1.3}
\begin{table}[h]
\centering
\begin{tabular}{|l|l|l|l|l|}
\hline
Dataset & Features & Size & Year & Organization \\ \hline
BDD100K & Driving scenes, multiple lighting conditions & 100GB & 2018 & Berkeley \\ \hline
\end{tabular}
\caption{Datasets}
\end{table}

% ---------------- RESULTS ----------------
\chapter{Results and Discussion}
\renewcommand{\arraystretch}{1.3}

\begin{tabular}{|l|l|}
\hline
Mode & Time \\ \hline
Sequential & 8 sec \\ \hline
Pipeline (Python) & 5 sec \\ \hline
OpenMP (C++) & 3 sec \\ \hline
\end{tabular}

% ---------------- REFERENCES ----------------
\chapter{References}

\begin{itemize}
\item OpenCV Documentation
\item YOLO Paper
\item BDD100K Dataset
\end{itemize}

\end{document}
